\section{Green potentials and the driving process}
\label{cal:section}

Throughout this section $m>0$.  We work under one of the subsequential
laws supplied by Proposition~\ref{geo:compactness}, in the joint filtration
of the curve and the fermion field of Proposition~\ref{geo:martingale}.
In particular, $M_t(w)$ is a continuous local martingale only while $w$
is in the active domain.  No martingale property after swallowing will
be used.  Proposition~\ref{eng:energy} supplies the other probabilistic input:
\begin{equation}
 \mathbb E\int_0^\infty J_t^2\,\dd t<\infty,
 \qquad
 J_t=\int_{\Hh}\frac{\mu_t(v)}{|v|^2}\,\dd A(v).
 \label{cal:energy-input}
\end{equation}
Here and below an expression involving an active point is used only
before its lifetime.  In a nonnegative space--time integral it is set
to zero after that lifetime.

It is convenient to abbreviate the physical uniformizing map by
\[
 \Phi_t=g_t\circ\phi,
 \qquad Z_t(w)=\Phi_t(w)-W_t,
 \qquad P_t(w)=P(Z_t(w)).
\]
Thus $\Phi_t$ and $Z_t$ map the physical active domain to $\Hh$,
while $\psi_t(v)=\Phi_t^{-1}(W_t+v)$ is the inverse centered chart.
The disc maps $\chi$ are used to compare complete curves; throughout
this section $t$ remains capacity time. We use the Dirichlet Green kernel
for $-\Delta$:
\begin{equation}
 G_{\Hh}(z,v)=\frac1{2\pi}
       \log\left|\frac{z-\overline v}{z-v}\right|,
 \qquad
 G_t(z,w)=G_{\Hh}(\Phi_t(z),\Phi_t(w)).
 \label{cal:green-normalization}
\end{equation}
The kernel $G_t$ is set to zero when either argument has left the active
domain.  The following proposition is the conclusion of the section.

\begin{proposition}[Identification of the driver]
\label{cal:driver}
The subsequential driving process is a continuous semimartingale and,
in the above joint filtration, satisfies
\begin{equation}
 \dd W_t=\sqrt{\frac{16}{3}}\,\dd B_t+
                \frac{2\pi}{3}C_t\,\dd t,
 \qquad
 \int_0^T C_t^2\,\dd t<\infty\quad\text{almost surely for every }T<\infty,
 \label{cal:driver-equation}
\end{equation}
where $B$ is a Brownian motion and
\begin{equation}
 C_t=\frac4\pi\int_{\Hh}P(v)\mu_t(v)
              |\nabla(h_t\circ\psi_t)(v)|\,\dd A(v).
 \label{cal:C-definition}
\end{equation}
In particular, neither absolute continuity of the quadratic variation
nor absolute continuity of the finite-variation part of $W$ is an
assumption of the argument.
\end{proposition}

We first express $h_t$ as harmonic measure plus a bounded Green
potential. A height cutoff makes this potential a semimartingale up to
swallowing; differences at two protected points then recover
semimartingality of $W$. A common covariance clock records possible
singular parts of both the driver and the fermion brackets. Finally,
averaged expansions at the moving mark identify the cubic and quadratic
poles, giving the variance, eliminating singular driver variation, and
determining the drift.

\subsection{The bounded potential and interior estimates}
\label{cal:deterministic}

With $D=2\ii\partial_z$, the relation $Dh_t=M_t^2$ reads
\[
 Dh=h_y+\ii h_x,
 \qquad |\nabla h_t|=|M_t|^2,
 \qquad \overline\partial M_t=-\ii m\overline{M_t}.
\]
Consequently $\Delta h_t=-4m|M_t|^2$.  With
$u_t(w)=\arg Z_t(w)/\pi$, the boundary values of $h_t-u_t$ are zero at
all prime ends other than the two marks.

\begin{lemma}[Potential representation and bulk bound]
\label{cal:potential}
There is a constant $K$, depending only on $m$ and a physical diameter
bound, such that on every active domain
\begin{align}
 h_t(z)&=u_t(z)+4m\int_{O_t}G_t(z,w)|M_t(w)|^2\,\dd A(w),
 \label{cal:representation}\\
 |M_t(w)|^2&\leq \frac{K}{\operatorname{crad}_{O_t}(w)}.
 \label{cal:bulk-bound}
\end{align}
In particular the nonnegative potential
\begin{equation}
 Q_t(z):=\int_{O_t}G_t(z,w)|M_t(w)|^2\,\dd A(w)
 \quad\text{satisfies}\quad 0\leq Q_t(z)\leq\frac1{4m}.
 \label{cal:bounded-potential}
\end{equation}
\end{lemma}

\begin{proof}
We first justify the gradient estimate for a bounded solution of
$\Delta h=-4m|\nabla h|$ in a bounded domain.  The following elementary
interior argument keeps its constants uniform as the domain is slit.
On a disc $B(x,R)$ compactly contained in the domain, let
\[
 N=\sup_{y\in B(x,R)}(R-|y-x|)|\nabla h(y)|.
\]
One can first replace the disc by a slightly smaller disc, so that
$N<\infty$.  Choose $y$ attaining at least $N/2$, and put
$d=R-|y-x|$.  On $B(y,d/2)$ the gradient is at most $2N/d$.
Decompose $h$ there into its harmonic extension of the boundary data
and its Green potential.  The derivative of the Poisson kernel at the
centre, and the integrability of the derivative of the disc Green
kernel, give
\[
 |\nabla h(y)|\leq \frac{K}{d}\|h\|_\infty+
       Kd\,\|\Delta h\|_{L^\infty(B(y,d/2))}
 \leq \frac{K}{d}\|h\|_\infty+KmN.
\]
Multiplication by $d$ absorbs the last term when $mR$ is below a fixed
constant.  Taking
$R$ comparable to $\min\{\operatorname{dist}(x,\partial O_t),m^{-1}\}$,
and using $0\leq h_t\leq1$, proves
\[
 |\nabla h_t(x)|\leq
 K\bigl(\operatorname{dist}(x,\partial O_t)^{-1}+m\bigr).
\]
The distance is bounded above by the physical diameter.  Koebe's
one-quarter estimate compares distance with conformal radius and gives
\eqref{cal:bulk-bound}.  Interior regularity needed in this argument
also follows directly from the constant-coefficient equation for $M$;
indeed $\Delta M=4m^2M$.

For the representation, fix $t$ and write $k=h_t-u_t$.  Exhaust $O_t$
by images of concentric discs of radii $r\uparrow1$ under an interior
normalized Riemann map.  Green's formula in each exhaustion domain is
\[
 k(z)=H_r(z)+4m\int_{O_t^{(r)}}
                   G_{O_t^{(r)}}(z,w)|M_t(w)|^2\,\dd A(w),
\]
where $H_r$ is the harmonic extension of $k$ from the exhaustion
boundary.  In disc coordinates the functions $k$ are uniformly bounded
and tend to zero radially at every boundary point except the two
marks.  The Poisson formula and dominated convergence imply
$H_r(z)\to0$.  The Green kernels increase to $G_t$; monotone convergence
therefore proves \eqref{cal:representation}, as well as nonnegativity
of $k$.  Since $0\leq h_t,u_t\leq1$, this also proves
\eqref{cal:bounded-potential}.  This reasoning uses prime-end boundary
values and requires no physical boundary regularity.
\end{proof}

At a fixed time, put
\[
 F_t(v)=\psi_t'(v)^{1/2}M_t(\psi_t(v)),
\]
using either consistent square-root choice.  Conformal radius and area
transform as
\begin{equation}
 \operatorname{crad}_{O_t}(\psi_t(v))=2\operatorname{Im}v\,|\psi_t'(v)|,
 \qquad
 m|M_t(\psi_t(v))|^2\,\dd A(\psi_t(v))
       =\mu_t(v)|F_t(v)|^2\,\dd A(v).
 \label{cal:changes-of-variables}
\end{equation}
It follows that
\begin{equation}
 |F_t(v)|^2\leq\frac K{\operatorname{Im}v},
 \qquad
 m\int_{O_t}P_t(w)|M_t(w)|^2\,\dd A(w)\leq KJ_t,
 \qquad C_t\leq KJ_t.
 \label{cal:poisson-control}
\end{equation}
For the last inequality we used
$|\nabla(h_t\circ\psi_t)|=|F_t|^2$.
Also, bounded physical area gives the deterministic identity
\begin{equation}
 \int_{\Hh}\mu_t(v)^2\,\dd A(v)=m^2|O_t|<\infty.
 \label{cal:mu-L2}
\end{equation}

\subsection{Differentiating the potential without a killing term}
\label{cal:cutoff-subsection}

A direct differentiation of the logarithmic kernel, using the Loewner
ordinary differential equation, gives
\begin{equation}
 \partial_tG_t(z,w)=-\frac2\pi P_t(z)P_t(w).
 \label{cal:hadamard}
\end{equation}
This identity differentiates only $\Phi_t(z)$ and $\Phi_t(w)$.
In particular it holds for every continuous driver and does not assume
that $W$ is a semimartingale.

We will localize at a protected physical point $z$: the time horizon is
bounded, $z$ has a fixed open neighbourhood in the active domain, and
its chart height, coordinates, and derivative stay in compact ranges
with positive lower bounds where appropriate.  These are adapted
stopping localizations.  A countable family of physical discs and
positive rational protection levels suffices.  The essential
protection is a positive lower bound on chart height; positive physical
distance alone would not exclude swallowing inside a disconnected
pocket.

\begin{lemma}[Stochastic Green potential]
\label{cal:stochastic-potential}
On every such stopped interval $[0,T]$,
\begin{equation}
 \mathbb E\int_0^T\int_{O_t}G_t(z,w)
                         \,\dd[M(w),\overline{M(w)}]_t\,\dd A(w)<\infty.
 \label{cal:weighted-bracket}
\end{equation}
The potential $Q_t(z)$ has a continuous semimartingale version, and
\begin{align}
 \dd Q_t(z)={}&-\frac2\pi P_t(z)
       \left(\int_{O_t}P_t(w)|M_t(w)|^2\,\dd A(w)\right)\dd t
 \notag\\
 &+\int_{O_t}G_t(z,w)\,\dd[M(w),\overline{M(w)}]_t\,\dd A(w)
 \notag\\
 &+2\operatorname{Re}\int_{O_t}
                G_t(z,w)\overline{M_t(w)}\,\dd M_t(w)\,\dd A(w).
 \label{cal:potential-differential}
\end{align}
The last line is defined by the weighted bracket-isometry limit of
compactly supported stochastic integrals.  The statement holds
simultaneously for countably many protected spatial points, and also
as a spatial distribution identity on protected open sets.
\end{lemma}

\begin{proof}
Choose a smooth nondecreasing function $\chi:[0,\infty)\to[0,1]$ which
is zero on $[0,1]$ and one on $[2,\infty)$.  Write
$Y_t(w)=\operatorname{Im}\Phi_t(w)$, and insert the factor
$\chi(Y_t(w)/\lambda)$ in $Q_t$.  First also insert time-independent
compact physical cutoffs $0\leq\rho_N\leq1$ increasing to one.  Denote
the resulting potential by $Q_t^{N,\lambda}(z)$.

For fixed $N,\lambda$ the product formula is legitimate entirely before
lifetimes.  Here is one explicit way to implement this assertion.
For each $w$ in the compact support, stop its martingale when
$Y_t(w)\leq\lambda$, and freeze it thereafter.  On the set where the
cutoff is nonzero this stopped martingale equals $M_t(w)$.  Moreover,
on a bounded time horizon,
\[
 |\Phi_t'(w)|\leq |\phi'(w)|\exp(2T/\lambda^2)
       \quad\text{as long as }Y_t(w)\geq\lambda.
\]
Together with \eqref{cal:bulk-bound}, this bounds these stopped
martingales uniformly over the compact support.  Their pointwise
expected brackets are therefore uniformly bounded.  The kernel is
bounded above by $G_\Omega(z,w)$; its only interior singularity is
logarithmic and is integrable in $w$.  These bounds justify stochastic
Fubini for the compact integral.  Equivalently, approximate its spatial
integrals by simple functions and use the covariance Cauchy--Schwarz
inequality.  The singularity can also be tested first against a smooth
compact function of $z$.

Both time-dependent factors multiplying $|M|^2$ decrease.  Formula
\eqref{cal:hadamard} gives one nonnegative loss.  The other is
\[
 -\dd\chi(Y_t/\lambda)
   =\frac{2Y_t}{\lambda|Z_t|^2}\chi'(Y_t/\lambda)\,\dd t.
\]
Take $\lambda$ smaller than a fixed fraction of the protected height
of $z$.  On $\lambda\leq Y_t(w)\leq2\lambda$ the explicit half-plane
kernel gives $G_t(z,w)\leq K_zY_t(w)$.  Including the area Jacobian
from \eqref{cal:changes-of-variables}, the absolute cutoff loss is at
most
\begin{align}
 &K_z\int_{\lambda\leq\operatorname{Im}v\leq2\lambda}
       \frac{(\operatorname{Im}v)^2}{\lambda|v|^2}
       \frac{\mu_t(v)}m|F_t(v)|^2\,\dd A(v)\,\dd t
 \notag\\
 &\hspace{25mm}\leq
 K_z\int_{\operatorname{Im}v\leq2\lambda}
                      \frac{\mu_t(v)}{|v|^2}\,\dd A(v)\,\dd t.
 \label{cal:cutoff-loss}
\end{align}
Constants may depend on the fixed mass and the protection of $z$, but
not on $N$.  By \eqref{cal:energy-input},
$\mathbb E\int_0^T J_t\,\dd t<\infty$.  Dominated convergence shows
that the expected total loss in \eqref{cal:cutoff-loss} tends to zero
as $\lambda\downarrow0$.  The ordinary kernel loss is bounded by
$K_zJ_t\,\dd t$, by \eqref{cal:poisson-control}.

Apply the product formula to $Q^{N,\lambda}$ and take expectations.
Its bracket contribution is nonnegative, and its initial and terminal
potentials lie between zero and $1/(4m)$.  Consequently
\[
 \mathbb E\int_0^T\int
 \rho_N(w)\chi(Y_t(w)/\lambda)G_t(z,w)
                \,\dd[M(w),\overline{M(w)}]_t\,\dd A(w)
 \leq K_z\left(1+\mathbb E\int_0^T J_t\,\dd t\right).
\]
If an additional martingale localization is used in deriving this
inequality, positivity allows its removal by Fatou's lemma.
Letting $N\uparrow\infty$ and $\lambda\downarrow0$ monotonically
proves \eqref{cal:weighted-bracket}.

For completeness, the covariance inequality controlling the
martingale integrals is
\begin{equation}
 \begin{aligned}
 &\left|\int G_t(z,w)\,2\operatorname{Re}
                         (\overline{M_t(w)}A_t(w))\,\dd A(w)\right|^2\\
 &\qquad\leq 4\left(\int G_t(z,w)|M_t(w)|^2\,\dd A(w)\right)
         \left(\int G_t(z,w)|A_t(w)|^2\,\dd A(w)\right).
 \end{aligned}
 \label{cal:weighted-isometry}
\end{equation}
Here $A$ can be any Hilbert realization of a finite family of
covariance densities; approximation gives the same inequality for
spatial integrals.  For the difference of two cutoffs one inserts the
absolute cutoff difference in both factors.  The first factor is
uniformly bounded by \eqref{cal:bounded-potential}, and the expected
integral of the second tends to zero by
\eqref{cal:weighted-bracket}.  Thus the martingale terms converge in
$L^2$, uniformly in probability on the stopped time interval.  The
positive bracket terms converge in expected total variation.  The
kernel-loss terms converge in $L^1$, and the cutoff loss disappears
by \eqref{cal:cutoff-loss}.  At each fixed time the potentials converge
by their nonnegative integral representation.  Passing to the limit
in the product identity proves \eqref{cal:potential-differential} and
its continuous version.  In particular the cutoff does not leave a
measure concentrated on swallowing times.  The positive bracket
limits are continuous as well: along a subsequence their total
variation tails converge uniformly to zero, and the approximating
bracket integrals are continuous.

The construction applies to every protected point in a countable
dense set.  Spatially testing before applying the same argument gives
the distributional statement.  Local elliptic estimates, made
explicit in Lemma~\ref{cal:covariance}, identify these versions with
point evaluations wherever those are used below.
\end{proof}

\subsection{Recovering semimartingality of the driver}
\label{cal:recovery-subsection}

\begin{lemma}[Recovery and global localization]
\label{cal:recovery}
The continuous process $W$ is a semimartingale in the joint filtration.
\end{lemma}

\begin{proof}
Take $z,z'$ joined by a fixed path inside a protected physical disc.
The identity
\[
 h_t(z)-h_t(z')=\operatorname{Im}\int_{z'}^z M_t(\zeta)^2\,\dd\zeta
\]
shows that this difference is a continuous semimartingale.  Indeed the
field is uniformly bounded after protection, the expected brackets
are uniformly bounded on the compact path, and covariance
Cauchy--Schwarz justifies the one-dimensional spatial stochastic
integral in its product formula.  Subtracting
\eqref{cal:potential-differential} at the two points proves that
$u_t(z)-u_t(z')$ is a semimartingale.

For complex $a,a'\in\Hh$ and real $w$, define
\[
 H(a,a',w)=\frac1\pi\bigl(\arg(a-w)-\arg(a'-w)\bigr).
\]
Its derivative in the real variable is
\[
 \partial_wH(a,a',w)=\frac1\pi\bigl(P(a-w)-P(a'-w)\bigr).
\]
Whenever $P_t(z)\ne P_t(z')$, the implicit-function theorem therefore
recovers $W_t$ smoothly from
$u_t(z)-u_t(z')$, $\Phi_t(z)$, and $\Phi_t(z')$ in a neighbourhood of
the current state.  The last two processes have finite variation by
the Loewner equation.  This proves semimartingality on that
neighbourhood without applying stochastic calculus to $W$ in advance.

We give a stopping localization that covers every finite time, so
that repeated swallowings cannot leave an exceptional accumulation
time.  Enumerate rational physical discs with closure in $\Omega$,
rational positive chart-protection levels, and finite collections of
rational pairs in each disc.  Stop before chart height on the closed
disc falls below its specified level, before the coordinates and
$W$ exceed specified bounds, or at a specified finite horizon.  All
these quantities are continuous before loss of protection.  Also stop
when the maximum of the absolute Poisson-kernel differences over the
finite collection falls below a specified positive threshold.
If the initial inequalities fail, set the stopping time equal to zero.

For each such choice, the allowed coordinate states form a compact
set on which at least one of the finitely many derivatives
$\partial_wH$ is bounded away from zero.  Cover this compact set by a
finite collection of implicit-function neighbourhoods, each with a
smaller neighbourhood whose closure is contained in the first.
Select a valid smaller neighbourhood measurably, and change charts
only on leaving the corresponding larger neighbourhood.  The fixed
positive separation and continuity of the finitely many coordinate
processes prevent an accumulation of chart changes on a protected
finite interval.  For clarity, let $S_k$ be the successive chart-change times, capped at
this protection stop $\sigma$ and kept equal to $\sigma$ once it is
reached.  Each finite concatenation gives the semimartingale
$W^{S_k}$; the chart selected at $S_k$ is $\mathcal F_{S_k}$-measurable,
so its stopped increment is a finite sum of predictable stopped
integrals of the corresponding smooth implicit function.  Define
\[
 R_k=\begin{cases}S_k,&S_k<\sigma,\\
                  \infty,&S_k=\sigma.
       \end{cases}
\]
These are increasing stopping times, and the nonaccumulation of chart
changes gives $R_k\uparrow\infty$.  Since
$(W^\sigma)^{R_k}=W^{S_k}$, they are ordinary semimartingale localizers
for $W^\sigma$.  Thus $W$ stopped at every enumerated protection time
is a semimartingale.

To see that these times cover all horizons, fix a sample path and
$T<\infty$.  The active domain $O_T$ is nonempty; choose a rational
closed disc contained in it.  This disc survives on $[0,T]$, and
compactness provides positive chart-height and finite coordinate
bounds there.  The function $P_t$ is not constant on this disc:
\[
 |\nabla P_t(z)|=\frac{|\Phi_t'(z)|}{|Z_t(z)|^2}>0.
\]
At each time some rational pair has different Poisson-kernel values.
Continuity and compactness of $[0,T]$ select finitely many pairs whose
maximum difference stays above a positive threshold.  Some
localization in the enumeration therefore lasts beyond $T$.

Finally, finite maxima of semimartingale stopping localizations are
again such localizations, using
$W^{\sigma\vee\tau}=W^\sigma+W^\tau-W^{\sigma\wedge\tau}$.
The successive finite maxima tend to infinity almost surely, by the
previous paragraph.  This is a global semimartingale localization
sequence.  The normalized continuous local-martingale and continuous
finite-variation decompositions agree on overlaps, because a continuous
finite-variation local martingale is constant.  They therefore glue
to the canonical decomposition
\begin{equation}
 W=W^{\mathrm M}+V,
 \label{cal:decomposition}
\end{equation}
where $W^{\mathrm M}$ is a continuous local martingale and $V$ has locally finite
variation. At this stage $V$ and $[W^{\mathrm M}]$ may have singular parts.
\end{proof}

\subsection{A simultaneous covariance clock}
\label{cal:clock-subsection}

The next construction is needed because the driver decomposition in
\eqref{cal:decomposition} may still charge sets of zero capacity time.
Pointwise covariance densities chosen separately at each spatial point
would not be sufficient: later we will choose a chart and balls
approaching its moving boundary mark after fixing a time. The clock $L$
below serves this simultaneous covariance construction; it differs from
the point-dependent clock $\sigma$ used for geometric estimates in
Section~\ref{eng:section}.

\begin{lemma}[Spatial covariance realization]
\label{cal:covariance}
Locally in time there is a continuous increasing process $L$ such that
$\dd t$, $|\dd V|$, $\dd[W^{\mathrm M}]$, and all the spatial bracket measures
used below are absolutely continuous with respect to $\dd L$.  Outside one null set for $\dd L\,\dd\mathbb P$, write
\begin{equation}
 \tau=\frac{\dd t}{\dd L},\qquad
 \nu=\frac{\dd V}{\dd L},\qquad
 U=\dot W^{\mathrm M},\qquad A_{\mathrm{phys}}(w)=\dot M(w).
 \label{cal:clock-notation}
\end{equation}
The dots denote covariance vectors with respect to $\dd L$, not time
derivatives. These vectors lie in a real Hilbert space, complexified
for $A_{\mathrm{phys}}$.  Their scalar products are the corresponding
bracket densities.  The product is extended bilinearly, so
$A\cdot\overline A=|A|^2$ and $U\cdot U=|U|^2$.
The field $A_{\mathrm{phys}}$ is spatially smooth on every compact
subset of the live domain.  Spatial evaluations, derivatives, compact
tests, and the Green stochastic integrals in
\eqref{cal:potential-differential} have their covariance vectors on this
same full set.  In particular, for protected $z$,
\begin{equation}
 \int_{O_t}G_t(z,w)|A_{\mathrm{phys}}(w)|^2\,\dd A(w)<\infty
 \quad\text{for }\dd L\text{-almost every live time.}
 \label{cal:clock-green-integrability}
\end{equation}
\end{lemma}

\begin{proof}
We describe the construction on a countable exhaustion; this also
specifies the simultaneous exceptional set.  Let $B$ range over
rational physical discs with closure in $\Omega$.  Stop the entire
field on $B$ at one common stopping time, before the minimum chart
height on its closure falls below a positive rational level and before
a fixed time or coordinate bound is exceeded.  Include slightly larger
discs whenever interior evaluations on $B$ will be used.  The resulting
stopped field is uniformly bounded, by
\eqref{cal:bulk-bound} and the Loewner derivative estimate in the proof
of Lemma~\ref{cal:stochastic-potential}.  Freezing the whole field at
the same time preserves its equation in the spatial variable.

Let $\mathcal H_B$ be the real Hilbert space of square-integrable
complex functions on $B$ which solve
$\overline\partial f+\ii m\overline f=0$ distributionally, with inner
product $\operatorname{Re}\int_B f\overline g\,\dd A$.
It is a closed subspace of real $L^2(B;\mathbb C)$ and hence is
separable.  Interior estimates for $\Delta f=4m^2f$ show that evaluation
and every spatial derivative on a smaller disc are bounded linear
functionals on $\mathcal H_B$.  Their Riesz representatives depend
smoothly on the evaluation point: the Taylor remainder estimates
follow from the same interior bounds one derivative further.

The stopped field is a continuous square-integrable
$\mathcal H_B$-valued martingale.  To verify this assertion, integrate
its bounded scalar martingale identities against a countable dense set
of $L^2$ tests; approximation gives all Hilbert-space conditional
expectation identities.  Joint local continuity gives Hilbert
continuity.  If $(e_j)$ is a real orthonormal basis, its trace bracket is
\[
 \sum_j[\langle M,e_j\rangle]_t,
 \qquad
 \mathbb E\sum_j[\langle M,e_j\rangle]_T
   =\mathbb E\|M_T-M_0\|_{\mathcal H_B}^2<\infty.
\]
Enumerate all discs and protections.  Multiplying their centred
stopped martingales by sufficiently small positive deterministic
weights makes their direct sum a continuous square-integrable
martingale in a separable real Hilbert space.  This is possible because
each stopped field has a deterministic norm bound.  For the driver use the normalized martingale
\[
 N_t^0=\int_0^t\frac{\dd W_s^{\mathrm M}}{1+[W^{\mathrm M}]_s},
 \qquad [N^0]_t=\frac{[W^{\mathrm M}]_t}{1+[W^{\mathrm M}]_t}\leq1.
\]
It is square-integrable and can be included as one real coordinate.
The joint covariance of this direct sum retains all cross-brackets,
including those between different discs and this normalized driver.
The covariance vector of the original driver is recovered by
multiplying that coordinate's vector by $1+[W^{\mathrm M}]_t$.

A positive trace-class covariance measure has an operator density
with respect to its trace measure.  This can be constructed directly:
take Radon--Nikodym derivatives of its matrix entries in a countable
orthonormal basis, impose positivity on the countable collection of
rational finite linear combinations, and complete the resulting
positive form.  Its finite trace makes it a bounded positive
trace-class operator at almost every clock time.  Its positive square
root applied to the Riesz representatives of evaluation gives the
vectors in \eqref{cal:clock-notation}.  Real and imaginary evaluations
supply the complexified vector $A_{\mathrm{phys}}$.
Smoothness of the Riesz representatives proves smoothness in $w$.
On overlapping protected discs the evaluation martingales coincide
before either stop; the bracket of their difference is zero there.
Impose this equality first at rational spatial points and then use
smoothness.  The local realizations therefore agree on their overlaps
on one common full set.

Add $\dd t$, $|\dd V|$, and $\dd[W^{\mathrm M}]$ to the trace clock.  Also add
the countably many positive Green-integrated bracket measures from
\eqref{cal:weighted-bracket}, using rational protected points or
nonnegative compact tests in $z$.  Their cumulative processes are
finite on their stopped intervals.  There is no issue in making this
sum locally finite: for any continuous increasing process $H$, the
measure $(1+H_t)^{-2}\dd H_t$ has total mass at most one and has the
same null sets as $\dd H$; use these measures with summable positive
weights.  Every required measure is therefore absolutely continuous with respect
to $\dd L$.  Matrix
Radon--Nikodym derivatives with respect to this clock give the
same construction, now including singular parts of $\dd V$ or
$\dd[W^{\mathrm M}]$.

Choose the countable smooth test families dense in the $C^\infty$
topology for each fixed compact support in a countable protected
exhaustion, and include a nonzero nonnegative bump in each protected
disc.  For each such nonnegative test, nonnegative Fubini identifies
the density of its Green-integrated bracket measure with the
corresponding integral of $G_t|A_{\mathrm{phys}}|^2$.  Fix a time in
the resulting common full set and a surviving bump $\eta$, and put
\[
 K_\eta(w)=\int\eta(z)G_t(z,w)\,\dd A(z).
\]
Then $\int K_\eta|A_{\mathrm{phys}}|^2\,\dd A<\infty$.
For $z$ in any compact subset of the live domain, the explicit
half-plane kernel bounds $G_t(z,w)$ and its first $z$-derivatives by
a constant times $K_\eta(w)$ outside a slightly larger compact subset.
Locally, their logarithmic and inverse-distance singularities are
integrable against the smooth spatial fields.  This proves
\eqref{cal:clock-green-integrability} at every live interior point
on the same full set.  Weighted Cauchy--Schwarz gives the corresponding
tail control for the mixed Green integrals; \eqref{cal:weighted-isometry}
and \eqref{cal:bounded-potential} identify their compact-truncation
limits with the covariance vectors of the noncompact martingale
integrals.  The potential expressions define continuous spatial
distributions, so density of the chosen tests extends the identities
and their distributional derivatives to all compact tests on this
same full set.  An arbitrary current conformal chart, and the
boundary-approaching balls chosen later, are deterministic spatial
choices after a time in this full set has been fixed.
\end{proof}

At such a time, translate the half-plane to put $W_t$ at zero and omit
the time subscript.  Set
\[
 A(v)=\psi_t'(v)^{1/2}A_{\mathrm{phys}}(\psi_t(v)),
 \qquad F(v)=\psi_t'(v)^{1/2}M_t(\psi_t(v)),
 \qquad \mu(v)=m|\psi_t'(v)|.
\]
The same square root is used in the two definitions.  Formula
\eqref{cal:clock-green-integrability} implies
\begin{equation}
 \int_{\Hh}\frac{\operatorname{Im}v}{1+|v|^2}
                    \mu(v)|A(v)|^2\,\dd A(v)<\infty.
 \label{cal:weighted-A}
\end{equation}
Indeed, for every fixed $z_0\in\Hh$ there is $c(z_0)>0$ such that
\[
 G_{\Hh}(z_0,v)\geq
       c(z_0)\frac{\operatorname{Im}v}{1+|v|^2},\qquad v\in\Hh.
\]
This follows from the explicit kernel near the real line and infinity,
and positivity on the remaining compact set; its pole only increases
the ratio.  A Green average against a protected nonnegative test has
the same lower bound.  Finally,
\[
 |A_{\mathrm{phys}}(\psi_t(v))|^2\,\dd A(\psi_t(v))
       =\frac{\mu(v)}m|A(v)|^2\,\dd A(v),
\]
which proves \eqref{cal:weighted-A} on the same full clock-measure set.

\subsection{Three identities in the current half-plane}
\label{cal:identities-subsection}

All the next identities are first identities of spatial distributions;
the estimates that follow also justify the displayed kernels directly.
Every occurrence of $DG$ differentiates the first variable.  We claim
\begin{align}
 F(z)^2={}&\frac1{\pi z}
       +4\int_{\Hh}DG(z,v)\mu(v)|F(v)|^2\,\dd A(v),
 \label{cal:identity-zero}\\
 2F(z)A(z)={}&\frac{U}{\pi z^2}
       +8\int_{\Hh}DG(z,v)\mu(v)
                    \operatorname{Re}(\overline{F(v)}A(v))\,\dd A(v),
 \label{cal:identity-martingale}\\
 A(z)\cdot A(z)={}&\frac{|U|^2-4\tau}{\pi z^3}
       +\frac\nu{\pi z^2}-\frac{2C\tau}{z^2}
       +4\int_{\Hh}DG(z,v)\mu(v)|A(v)|^2\,\dd A(v).
 \label{cal:identity-drift}
\end{align}
When $\tau=0$, the expression $C\tau$ is defined to be zero; it is a
capacity-density term and does not require a value for $C$ on that
part of the clock.

Here is a derivation including the constants.  In physical coordinates,
\[
 Du_t(w)=\frac{\Phi_t'(w)}{\pi Z_t(w)},\qquad
 \partial_t\Phi_t'(w)=-\frac{2\Phi_t'(w)}{Z_t(w)^2}.
\]
It\^o's formula, now justified by Lemma~\ref{cal:recovery}, gives
\begin{equation}
 \dd\left(\frac{\Phi_t'(w)}{\pi Z_t(w)}\right)
 =\frac{\Phi_t'(w)}{\pi Z_t(w)^2}\,\dd W_t
   +\frac{\Phi_t'(w)}{\pi Z_t(w)^3}
                    \bigl(\dd[W^{\mathrm M}]_t-4\dd t\bigr).
 \label{cal:ito-harmonic}
\end{equation}
The two contributions $-2\dd t$ in the last coefficient come from
$\dd(1/Z_t)$ and from $\dd\Phi_t'$, respectively.
Applying $D$ to \eqref{cal:representation} gives
\eqref{cal:identity-zero} after division by $\Phi_t'$.  Applying $D$
to \eqref{cal:potential-differential}, and comparing with
\[
 \dd(M_t^2)=2M_t\,\dd M_t+\dd[M,M]_t,
\]
gives the other two identities.  In particular the kernel-loss term
has derivative
\[
 -\frac{8m}{\pi}\,DP_t(w)
       \int_{O_t}P_t(\zeta)|M_t(\zeta)|^2\,\dd A(\zeta)\,\dd t.
\]
Since $DP(z)=z^{-2}$ and
\[
 C=\frac{4m}{\pi}\int_{O_t}P_t(\zeta)|M_t(\zeta)|^2\,\dd A(\zeta),
\]
this becomes $-2C\tau/z^2$ after division by $\Phi_t'$ and $\dd L$.
The martingale potential has factor $4m\times2=8m$.
Changing physical area to chart area turns
$m\operatorname{Re}(\overline M A_{\mathrm{phys}})\dd A$
into $\mu\operatorname{Re}(\overline F A)\dd A$, and turns
$m|A_{\mathrm{phys}}|^2\dd A$ into $\mu|A|^2\dd A$.
These are precisely the factors eight and four in
\eqref{cal:identity-martingale}--\eqref{cal:identity-drift}.
This procedure differentiates identities in physical space first and
only then changes coordinates at a fixed time; it does not apply
It\^o's formula to a fermion evaluated at a random moving point.

\subsection{Averaged estimates at the moving mark}
\label{cal:tip-subsection}

For $s>0$ let
\[
 B_s=B(\ii s,s/4),\qquad N_s=B(\ii s,s/2).
\]
For a scalar or Hilbert-valued function $q$, write
\[
 \|q\|_{L^p_{\rm av}(B_s)}
   =\left(\frac1{|B_s|}\int_{B_s}|q|^p\,\dd A\right)^{1/p}.
\]
Constants in this subsection may depend on the fixed time and its
finite clock densities.  Each little-oh assertion concerns
$s\downarrow0$ at that time.

\begin{lemma}[Tip expansions]
\label{cal:tip-expansions}
In the current conformal setting, assume
\eqref{cal:mu-L2}, \eqref{cal:weighted-A}, and the bulk bound
$|F(v)|^2\leq K/\operatorname{Im}v$ from
\eqref{cal:poisson-control}. Then
\eqref{cal:identity-zero}--\eqref{cal:identity-drift} imply
\begin{align}
 F(z)^2&=\frac1{\pi z}+o(s^{-1})
                   \qquad\text{uniformly on }B_s,
 \label{cal:tip-leading}\\
 2F(z)A(z)&=\frac U{\pi z^2}+o(s^{-3/2})
                   \qquad\text{in }L^2_{\rm av}(B_s),
 \label{cal:tip-mixed-weak}\\
 \int_{\Hh}DG(z,v)\mu(v)|A(v)|^2\,\dd A(v)
             &=o(s^{-2})
                   \qquad\text{in }L^1_{\rm av}(B_s).
 \label{cal:tip-bracket}
\end{align}
If in addition $J=\int_{\Hh}\mu(v)|v|^{-2}\,\dd A(v)<\infty$, then
\begin{align}
 F(z)^2&=\frac1{\pi z}+C+o(1)
                   \qquad\text{uniformly on }B_s,
 \label{cal:tip-constant}\\
 2F(z)A(z)&=\frac U{\pi z^2}+o(s^{-1})
                   \qquad\text{in }L^2_{\rm av}(B_s).
 \label{cal:tip-mixed-strong}
\end{align}
No finiteness of $J$ is required for
\eqref{cal:tip-leading}--\eqref{cal:tip-bracket}.
\end{lemma}

\begin{proof}
Differentiating \eqref{cal:green-normalization} gives
\begin{equation}
 DG(z,v)=\frac{\operatorname{Im}v}
                    {\pi(z-\overline v)(z-v)}.
 \label{cal:DG-explicit}
\end{equation}
For $z\in B_s$ and $v\notin N_s$ this implies
\begin{equation}
 |DG(z,v)|\leq K\frac{\operatorname{Im}v}{(s+|v|)^2}.
 \label{cal:kernel-far}
\end{equation}
For $v\in N_s$ the additional singular bound is
\begin{equation}
 |DG(z,v)|\leq \frac K{|z-v|}.
 \label{cal:kernel-near}
\end{equation}
The first bound follows by comparing both denominator factors to
$s+|v|$; when $|v|$ is comparable to $s$, the exclusion of $N_s$ keeps
$|z-v|$ at least $s/4$.  The second follows because
$|z-\overline v|\geq\operatorname{Im}z+\operatorname{Im}v$.
For each fixed $v\in\Hh$, uniformly as $z\in B_s$ tends to zero,
\begin{equation}
 DG(z,v)\longrightarrow \frac{P(v)}\pi.
 \label{cal:kernel-limit}
\end{equation}

We record two consequences of analyticity of $\psi'$.  Its subharmonic
mean inequalities on discs of radius comparable to $s$, contained in
$B(\ii s,3s/4)$, give
\begin{align}
 \sup_{N_s}\mu
 &\leq\frac K s
    \left(\int_{B(\ii s,3s/4)}\mu^2\,\dd A\right)^{1/2}
       =o(s^{-1}),
 \label{cal:mu-sup-L2}\\
 \sup_{N_s}\mu
 &\leq\frac K{s^2}\int_{B(\ii s,3s/4)}\mu\,\dd A
 \leq K\int_{B(\ii s,3s/4)}\frac\mu{|v|^2}\,\dd A=o(1)
       \quad\text{if }J<\infty.
 \label{cal:mu-sup-J}
\end{align}
The vanishing follows from absolute continuity of the corresponding
area integral.  Define the finite measure
\begin{equation}
 \dd\rho(v)=\frac{\operatorname{Im}v}{1+|v|^2}
                     \mu(v)|A(v)|^2\,\dd A(v).
 \label{cal:rho}
\end{equation}
On $N_s$, for small $s$, its density is comparable to
$s\mu|A|^2$.  In particular
\begin{equation}
 \int_{N_s}s\mu|A|^2\,\dd A=o(1).
 \label{cal:near-bracket-mass}
\end{equation}

We now estimate each potential separately.

\smallskip\noindent
\emph{The potential in \eqref{cal:identity-zero}.}
Off $N_s$, \eqref{cal:kernel-far} and $|F|^2\leq K/y$ bound it by
\[
 K\int_{\Hh}\frac\mu{(s+|v|)^2}\,\dd A=o(s^{-1}).
\]
To see the little-oh rather than merely a big-oh, split at a fixed
small radius $\varepsilon$.  On its inner side the $L^2$ norm of
$(s+|v|)^{-2}$ is $O(s^{-1})$, multiplied by an arbitrarily small
$L^2$ norm of $\mu$.  On the outer side the integral has a bound
independent of $s$, by $\mu\in L^2$.  First let $s$ tend to zero and
then let $\varepsilon$ tend to zero.  This also controls infinity.
On $N_s$, \eqref{cal:kernel-near}, $|F|^2\leq K/s$, and
$\int_{N_s}|z-v|^{-1}\dd A(v)\leq Ks$ give the uniform bound
$K\sup_{N_s}\mu=o(s^{-1})$.  This proves
\eqref{cal:tip-leading}.

If $J<\infty$, the far integrand is dominated by $K\mu/|v|^2$.
Equations \eqref{cal:kernel-limit} and dominated convergence give its
limit
\[
 \frac4\pi\int_{\Hh}P(v)\mu(v)|F(v)|^2\,\dd A(v)=C.
\]
The near part tends uniformly to zero by
\eqref{cal:mu-sup-J}.  The portion of the limiting integral removed
with $N_s$ tends to zero by the same $J$ majorant.  This proves
\eqref{cal:tip-constant}.

\smallskip\noindent
\emph{The mixed potential in \eqref{cal:identity-martingale}.}
Covariance Cauchy--Schwarz against the finite measure \eqref{cal:rho}
bounds the far part uniformly in $z\in B_s$ by
\begin{equation}
 K\rho(\Hh)^{1/2}
 \left(\int_{\Hh}\frac{(1+|v|^2)\mu(v)}{(s+|v|)^4}
                                             \,\dd A(v)\right)^{1/2}.
 \label{cal:mixed-far}
\end{equation}
Here the square of the kernel contributes $y^2$, and the weights
$|F|^2\leq K/y$ and $\dd\rho$ remove both factors of $y$.
The integral inside parentheses is $o(s^{-3})$: near zero this follows
from the $O(s^{-3})$ $L^2$ norm of $(s+|v|)^{-4}$ and the preceding
small-radius truncation argument; away from zero its kernel is
square-integrable, including its $|v|^{-2}$ tail.  Thus
\eqref{cal:mixed-far} is $o(s^{-3/2})$.
If $J<\infty$, multiply that integral by $s^2$ and write it against
$\mu(v)|v|^{-2}\dd A(v)$.  The remaining factor
\[
 \frac{s^2|v|^2(1+|v|^2)}{(s+|v|)^4}
\]
is bounded for $0<s\leq1$ and tends to zero for each fixed $v\ne0$.
Dominated convergence improves \eqref{cal:mixed-far} to $o(s^{-1})$.

For the near part put
$q(v)=\mu(v)\operatorname{Re}(\overline{F(v)}A(v))1_{N_s}(v)$.
Hilbert norms give
\begin{equation}
 \|q\|_{L^2(\Hh)}^2
 \leq \frac{K\sup_{N_s}\mu}{s^2}
                      \int_{N_s}s\mu|A|^2\,\dd A.
 \label{cal:mixed-near-input}
\end{equation}
The truncated kernel $1_{\{|v|<2s\}}/|v|$ has $L^1$ norm $O(s)$.
Young's inequality, followed by the factor $O(s^{-1})$ converting an
$L^2$ norm on $B_s$ to its normalized-area norm, therefore gives
\begin{equation}
 \begin{aligned}
 \|\text{mixed near part}\|_{L^2_{\rm av}(B_s)}
 &\leq \frac{K(\sup_{N_s}\mu)^{1/2}}s
       \left(\int_{N_s}s\mu|A|^2\,\dd A\right)^{1/2}\\
 &=o\bigl((\sup_{N_s}\mu)^{1/2}/s\bigr).
 \end{aligned}
 \label{cal:mixed-near}
\end{equation}
Hilbert-valued Young's inequality follows from the scalar convolution
inequality for the norm, so no finite-dimensional restriction is
involved.  Equations \eqref{cal:mu-sup-L2} and
\eqref{cal:mu-sup-J} yield precisely
\eqref{cal:tip-mixed-weak} and \eqref{cal:tip-mixed-strong}.
This step uses weighted integral control of $A$, not a pointwise
boundary bound.

\smallskip\noindent
\emph{The bracket potential in \eqref{cal:identity-drift}.}
Its far part is bounded by
\[
 K\int_{\Hh}\frac{1+|v|^2}{(s+|v|)^2}\,\dd\rho(v)=o(s^{-2}).
\]
Indeed multiplication of the integrand by $s^2$ makes it bounded
uniformly for $s\leq1$ and tending to zero pointwise.  Finite $\rho$
then gives dominated convergence.  For the near part, first integrate
the singular kernel in $z$:
\begin{align*}
 &\frac1{|B_s|}\int_{B_s}\int_{N_s}
          \frac{K\mu(v)|A(v)|^2}{|z-v|}\,\dd A(v)\,\dd A(z)\\
 &\hspace{20mm}\leq\frac K s\int_{N_s}\mu|A|^2\,\dd A
   =\frac K{s^2}\int_{N_s}s\mu|A|^2\,\dd A=o(s^{-2}).
\end{align*}
This proves \eqref{cal:tip-bracket} using only
\eqref{cal:weighted-A}, and completes the proof.
\end{proof}

\subsection{Pole comparison, including singular clock time}
\label{cal:comparison-subsection}

\begin{proof}[Proof of Proposition~\ref{cal:driver}]
Work at a time in the simultaneous full set of
Lemma~\ref{cal:covariance}.  The uniform leading expansion
\eqref{cal:tip-leading} implies
$|F(z)^2|\asymp s^{-1}$ on $B_s$ for all sufficiently small $s$.
Consequently division by $F^2$ is legitimate there.  Bilinearly squaring
\eqref{cal:tip-mixed-weak} and dividing by $4F^2$ gives
\begin{equation}
 A(z)\cdot A(z)=\frac{|U|^2}{4\pi z^3}+o(s^{-3})
                       \quad\text{in }L^1_{\rm av}(B_s).
 \label{cal:leading-square}
\end{equation}
For example, if $R=2FA-U/(\pi z^2)$, the cross-error is controlled by
$\|U/(\pi z^2)\|_{L^2_{\rm av}}\|R\|_{L^2_{\rm av}}$ and the
quadratic error by $\|R\|_{L^2_{\rm av}}^2$; multiplication by
$1/F^2=O(s)$ gives the displayed order.  This justifies the use of
averaged errors in a bilinear, rather than Hermitian, square.

Compare the cubic coefficients in
\eqref{cal:identity-drift}, using \eqref{cal:tip-bracket}.  This yields
\begin{equation}
 \frac{|U|^2}{4}=|U|^2-4\tau,
 \qquad |U|^2=\frac{16}{3}\tau.
 \label{cal:bracket-density}
\end{equation}
Coefficient comparison in these averaged statements can be read
without pointwise subsequences: set $z=s\zeta$, multiply by $s^3$,
and test on the fixed disc $B(\ii,1/4)$ against a smooth test with
nonzero pairing with $\zeta^{-3}$.  The lower-order and little-oh
terms vanish.

Consider first the singular part of the clock, where $\tau=0$.
Equation \eqref{cal:bracket-density} gives $U=0$.  Now the same weak
mixed estimate gives the stronger conclusion
\[
 A\cdot A=\frac{R\cdot R}{4F^2}=o(s^{-2})
                     \quad\text{in }L^1_{\rm av}(B_s).
\]
In \eqref{cal:identity-drift} only $\nu/(\pi z^2)$ can remain at
order $s^{-2}$, by \eqref{cal:tip-bracket}.  Rescaling and testing at
that order proves
\begin{equation}
 \nu=0\quad\text{on }\{\tau=0\}.
 \label{cal:singular-elimination}
\end{equation}
In particular this conclusion has not used $J<\infty$ on a set
singular to capacity time.

On $\{\tau>0\}$, \eqref{cal:energy-input} ensures $J<\infty$ at
$\dd L$-almost every time: a capacity-null set has zero
$\tau\dd L$ measure, hence zero $\dd L$ measure on $\{\tau>0\}$.
We can therefore use \eqref{cal:tip-constant} and
\eqref{cal:tip-mixed-strong}.  Uniformly on $B_s$,
\[
 \frac1{F^2}=\pi z-\pi^2Cz^2+o(s^2).
\]
The refined mixed remainder is $o(s^{-1})$ in normalized $L^2$;
its cross term, after division by $F^2$, is $o(s^{-2})$ in normalized
$L^1$.  We obtain
\begin{equation}
 A\cdot A=\frac{|U|^2}{4\pi z^3}
                   -\frac{|U|^2C}{4z^2}+o(s^{-2}).
 \label{cal:second-square}
\end{equation}
Comparing the quadratic coefficients in
\eqref{cal:identity-drift}, after cancelling the cubic ones, gives
\[
 \frac\nu\pi-2C\tau=-\frac{|U|^2C}{4},
 \qquad
 \nu=\frac{2\pi}{3}C\tau.
\]
Together with \eqref{cal:singular-elimination} this identity holds on
the whole clock.  Hence, in physical capacity time,
\[
 \dd[W^{\mathrm M}]_t=\frac{16}{3}\dd t,
 \qquad
 \dd V_t=\frac{2\pi}{3}C_t\,\dd t.
\]
The continuous martingale characterization of Brownian motion shows
that $B=\sqrt{3/16}\,W^{\mathrm M}$ is Brownian, with the harmless initial
constant chosen to make $W_0=0$.  Finally
\eqref{cal:poisson-control} and \eqref{cal:energy-input} give
\[
 \mathbb E\int_0^T C_t^2\,\dd t\leq
                       K\mathbb E\int_0^T J_t^2\,\dd t<\infty.
\]
This proves \eqref{cal:driver-equation}.  The stochastic proof used the
joint filtration only for the fermion martingales and their calculus.
The energy input itself concerns the marginal driver law, so its use
here does not require conditional crossing estimates in a larger,
possibly anticipative filtration.
\end{proof}
